\chapter{RESULTADOS E DISCUSSÕES}
\label{capitulo-resultados-discussoes}

    Ainda não tem nenhum resultado e discussão, pois o projeto ainda não foi concluído, marco necessário para confecção deste capítulo.
    
	Os resultados podem ser apresentados em forma de tabelas ou gráficos, sendo numerados sequencialmente e discutidos antes de serem colocados. Quando possível os resultados experimentais obtidos devem ser comparados com dados de literatura e suas diferenças (quando houver) discutidas.
	
	Compõe-se de: escrita do assunto a ser tratado no relatório, relato do fato, informação que se queira, descrição de norma e do procedimento. Apresentar somente assuntos que tenham relação direta e específica com o trabalho. 	Interprete as observações e os resultados obtidos. Procure estabelecer relação entre causa e efeito, partindo de suas observações. Mencione os problemas encontrados durante as atividades do projeto e procure relacionar os resultados obtidos com essas dificuldades. A discussão é a sua interpretação das observações e dos resultados. 
	
	Aqui o aluno deve apresentar os resultados esperados e pode também apresentar resultados semelhantes encontrados na literatura em estudos semelhantes realizados por outros autores que já pesquisaram a mesma área objeto da ação de extensão. 
	
	Caso tenha resultados prévios pode apresentar, mas nestes deve-se evitar a repetição por escrito de dados apresentados em tabelas e gráficos e as legendas desses devem ser autoexplicativas, dispensado a leitura do texto para compreensão. É muito importante explicar o que as respostas obtidas significam e como elas podem ajudar a solucionar o problema apresentado na extensão.
